\section*{Author conventions (publisher section 2)}
We denote by $\mathbb N$ the set of natural numbers and by $\mathbb N^+$ the set of positive natural numbers.
For every $n\in\mathbb N$, let $S_n$ be the symmetric group on $n$ elements. Let $\operatorname{Sym}(\mathbb N^+)$ be the symmetric group on $\mathbb N^+$.
For every $n\ge1$, $S_n$ can be embedded in $\operatorname{Sym}(\mathbb N^+)$ as the subgroup of all the $\rho\in\operatorname{Sym}(\mathbb N^+)$ such that $\rho i=i$ for all $i>n$. By abuse of notation, we shall write $S_n$ for the image of this embedding. We denote by $\iota$ the identity permutation in $\operatorname{Sym}(\mathbb N^+)$.
Given a semigroup $S$, we say that $u\in S$ is regular if there is $v\in S$ such that $uvu=u$. A semigroup is called regular if every element is regular.
Moreover, we say that $v$ is an inverse of $u$ if $uvu=u$ and $vuv=v$. Every element in a regular semigroup admits at least on inverse.
\paragraph{Definition 1.} An inverse semigroup is a semigroup $S$ endowed with a unary operation $(-)^*$ that satisfies, for all $u,v\in S$:
\[
(u^*)^*=u,\quad(uv)^*=v^*u^*,\quad uu^*u=u,\quad uu^*vv^*=vv^*uu^*.
\]
An inverse monoid is an inverse semigroup together with a unit for multiplication.
In an inverse semigroup $u^*$ is the unique inverse of $u$.
An idempotent in a semigroup $S$ is $u\in S$ such that $uu=u$. We denote by $E(S)$ the set of idempotents of $S$.
\paragraph{The natural order.} By definition, $u\le v$ if $uu^*v=u$, or, equivalently, if $vu^*u=u$. Actually, one can prove that the following are equivalent:
\begin{enumerate}
\item $u\le v$;
\item $u=ve$, for some idempotent $e$;
\item $u=dv$, for some idempotent $d$.
\end{enumerate}
Therefore, if $d$ is idempotent and $u\le d$, then $u$ is idempotent. Moreover, $(S,\le)$ is a partial order compatible with the structure of inverse semigroup, i.e., if $s\le t$ and $u\le v$, then $s^*\le t^*$ and $su\le tv$.
\paragraph{The minimum group congruence.} Two important relations in the study of an inverse semigroup $S$ are:
\begin{enumerate}
\item the relation of compatibility $\sim$, defined by $t\sim u$ if $t^*u$ and $tu^*$ are idempotents;
\item the relation $\sigma_S$, defined by $t\mathrel{\sigma_S}u$ if there exists $w$ such that $w\le t$ and $w\le u$.
\end{enumerate}
When $t$ and $u$ are compatible, i.e., $t\sim u$, then $t\land u$ exists and $t\land u=tu^*u$. The relation $\sigma_S$ is the smallest congruence such that $S/\sigma_S$ is a group [42, Thm.\ 2.4.1] and will be called minimum group congruence.
\paragraph{Definition 4.} An inverse semigroup $S$ is said to be
\begin{enumerate}
\item $E$-unitary if, whenever $e\in E(S)$ and $e\le v$, it follows that $v\in E(S)$;
\item $F$-inverse if each equivalence class of $\sigma_S$ has a unique maximal element.
\end{enumerate}
The inverse semigroup $S$ is $E$-unitary iff $\sim\,=\sigma_S$ [42, Thm.\ 2.4.6]; if $S$ is $F$-inverse, then it is $E$-unitary [42, Proposition 7.1.3].
\paragraph{Trees.} An ordered tree is a non-empty subset $A$ of $(\mathbb N^+)^*$ such that:
\begin{enumerate}
\item[(i)] $A$ is closed under initial segments;
\item[(ii)] $A$ is gapless: for every $i,j\in\mathbb N^+$ and every $w\in A$, if $i<j$ and $wj\in A$, then $wi\in A$.
\end{enumerate}
The elements of $A$ are called nodes; a node is a leaf if it has no children (immediate successors).
The empty word $\epsilon$ is the unique root, i.e., node without immediate predecessors.
We will be exclusively interested in finitely branching trees, i.e., trees whose nodes have a finite number of children.
We denote by $\operatorname{rd}(A)$ the degree of the root of $A$, i.e., the number of its children.
For each $1\le i\le\operatorname{rd}(A)$ we denote by $A_i:=\{w\in(\mathbb N^+)^*:iw\in A\}$ the subtree of $A$ of root $i$.
Finally, we say that a labelled tree is a function $l:A\to L$ from a tree to a nonempty set $L$ of labels. The tree $A$ is called the skeleton of $l$, and is denoted by $\operatorname{sk}(l)$. The set of finitely branching ordered trees can be defined coinductively. When we deal with infinite trees we will freely use coinductive arguments. We refer the reader to [38] for a gentle introduction to coinductive techniques.

\section*{Permutation trees (publisher section 3)}
\paragraph{Definition 12.} A permutation tree is a labelled ordered tree $t:A\to\operatorname{Sym}(\mathbb N^+)$ such that:
\begin{itemize}
\item if $a\in A$ has exactly $n\ge1$ children, then $t(a)\in S_n$;
\item if $a\in A$ is a leaf, then $t(a):=\iota$, where $\iota$ is the identity permutation in $\operatorname{Sym}(\mathbb N^+)$.
\end{itemize}
We let $\mathsf{PT}$ be the set of permutation trees.
We denote by $\bar\epsilon$ the unique permutation tree with skeleton $\{\epsilon\}$, i.e.\ the tree with one node labelled $\iota$.
We can conveniently write a permutation tree $t$ as $\langle\pi;t_1,\ldots,t_n\rangle$ where $\pi\in S_n$ is the label of the root, $\operatorname{rd}(t)=n$ and $t_1,\ldots,t_n\in\mathsf{PT}$.
Observe that $\mathsf{PT}$ can be defined coinductively as the greatest set $\mathsf{PT}$ such that, if $t\in\mathsf{PT}$, then there is $n\in\mathbb N$, $\pi\in S_n$ and $t_1,\ldots,t_n\in\mathsf{PT}$ such that $t=\langle\pi;t_1,\ldots,t_n\rangle$.
\paragraph{Definition 13.} Let $t,u\in\mathsf{PT}$. We define the product $tu$ by coinduction as follows. First, $t\bar\epsilon=t=\bar\epsilon t$.
If $t=\langle\pi;t_1,\ldots,t_n\rangle$, $u=\langle\rho;u_1,\ldots,u_m\rangle$ and $k=\max\{n,m\}$, then we define:
\[
tu:=\langle\pi\circ\rho;u_1t_{\rho1},\ldots,u_kt_{\rho k}\rangle
\]
where, by convention, $t_j=\bar\epsilon$ if $j>n$, and $u_i=\bar\epsilon$ if $i>m$.
The product operation makes $(\mathsf{PT},\cdot,\bar\epsilon)$ into a monoid. We now show that this monoid is an inverse monoid.
\paragraph{Definition 14.} Given $t\in\mathsf{PT}$, we define by coinduction $t^*$ as follows. First, $\bar\epsilon^*=\bar\epsilon$.
If $t=\langle\pi;t_1,\ldots,t_n\rangle$ with $n\ge1$, then $t^*=\langle\pi^{-1};(t_{\pi^{-1}1})^*,\ldots,(t_{\pi^{-1}n})^*\rangle$.
We denote by $E$ the set $\{t\in\mathsf{PT}:t(a)=\iota\text{ for every }a\in\operatorname{sk}(t)\}$.
Every $t\in E$ is univocally determined by its skeleton and we identify $\operatorname{sk}(t)$ with the permutation tree $t$.
\paragraph{Proposition 15.} The set $E(\mathsf{PT})$ of idempotent elements of $\mathsf{PT}$ coincides with $E$. Moreover, for every $d,e\in E$, $de=d\cup e$; consequently any two idempotents commute.
\begin{proof}
Let $t=\langle\pi;t_1,\ldots,t_n\rangle\in E(\mathsf{PT})$ be idempotent. Then $\pi\circ\pi=\pi$, and this is possible only if $\pi=\iota$.
The claim $t\in E$ follows from $t_i=(tt)_i=t_it_i$.
Conversely, if $d=\langle\iota;d_1,\ldots,d_n\rangle\in E$, then by definition of the product $dd=\langle\iota;d_1d_1,\ldots,d_nd_n\rangle$.
Since by coinduction $d_id_i=d_i$, we have $dd=d$ and $E=E(\mathsf{PT})$.
Finally, if $d,e\in E$, then, again by the coinductive definition of the product, $de=d\cup e$.
\end{proof}
\paragraph{Theorem 16.} The monoid $(\mathsf{PT},\cdot,(-)^*,\bar\epsilon)$ is an $E$-unitary inverse monoid.
\begin{proof}
We prove that $\mathsf{PT}$ is an inverse monoid by applying [31, Theorem 5.1.1], which states that $S$ is an inverse monoid iff $S$ is a regular monoid whose idempotents commute.
For every $t=\langle\pi;t_1,\ldots,t_n\rangle\in\mathsf{PT}$, we prove that $tt^*t=t$: $\pi\circ\pi^{-1}\circ\pi=\pi$ and, coinductively,
\[
(tt^*t)_i=((tt^*)t)_i=t_i(tt^*)_{\pi i}=t_i(t^*)_{\pi i}t_{(\pi^{-1}\circ\pi)i}=t_i(t_i)^*t_i=t_i.
\]
Since for every $t$ there is $t^*$ such that $tt^*t=t$, $\mathsf{PT}$ is regular. We conclude that $\mathsf{PT}$ is an inverse monoid applying Proposition 15.
We now prove that $\mathsf{PT}$ is $E$-unitary.
Let $t=\langle\pi;t_1,\ldots,t_n\rangle$, and $e=\langle\iota;e_1,\ldots,e_m\rangle$ be an idempotent such that $e\le t$.
If $e=\bar\epsilon$, then necessarily $t=\bar\epsilon$, hence $t\in E$.
If $e\ne\bar\epsilon$, then $e=td$ for some $d\in E$. We prove that $t\in E$.
First, we immediately get that $\pi=\iota$ and $n\le m$.
Moreover, for every $1\le i\le n$, we have $e_i=d_it_i$.
Since $d_i$ is idempotent, we get that $e_i\le t_i$, and coinductively, that $t_i\in E$. We conclude that $t\in E$.
\end{proof}
As a consequence of $E$-unitarity, if $t\le u$ and one between $t$ and $u$ is idempotent, then both are idempotent.
\paragraph{Lemma 18.} Let $t=\langle\pi;t_1,\ldots,t_n\rangle$ and $u=\langle\rho;u_1,\ldots,u_m\rangle$. Then $t\le u$ iff
\begin{enumerate}
\item $\pi=\rho$ and $m\le n$;
\item $t_i\le u_i$ for $1\le i\le m$ and $t_i\in E$ for $m<i\le n$.
\end{enumerate}
\begin{proof}
Assume that $t\le u$, meaning that $t=ue$ for some $e\in E$. Then $\pi=\rho$ and $n=\max\{m,\operatorname{rd}(e)\}\ge m$.
Moreover, since $t_i=e_iu_i$, for $1\le i\le m$, and $t_i=e_i$, for $m<i\le n$, the conclusion easily follows.
Conversely, let $e=\langle\iota;e_1,\ldots,e_m,t_{m+1},\ldots,t_n\rangle$, where $e_i$ is such that $t_i=e_iu_i$ with $e_i\in E$ for $1\le i\le m$. Then we have $t=ue$.
\end{proof}
Remark that the natural order is the greatest relation on $\mathsf{PT}$ that satisfies conditions (1) and (2) of Lemma 18.
\paragraph{Definition 19.} Given $t=\langle\pi;t_1,\ldots,t_n\rangle\in\mathsf{PT}$, we define coinductively $t^m$ as follows:
\[
t^m:=\begin{cases}
(\langle\pi;t_1,\ldots,t_{n-1}\rangle)^m,&\text{if }t_n\in E\text{ and }\pi n=n,\\
\langle\pi;(t_1)^m,\ldots,(t_n)^m\rangle,&\text{otherwise.}
\end{cases}
\]
Observe the similarity with the definition of Nakajima tree of a $\lambda$-term given in [7, Definition 2.60]. See also [6, Proposition 10.2.15].
Remark that if $e\in E$, then $e^m=\bar\epsilon$.
For $t\in\mathsf{PT}$, let $t^{\uparrow}:=\{u:t\le u\}$.
\paragraph{Proposition 20.} For every $t,u\in\mathsf{PT}$, the following hold:
\begin{enumerate}\item $t\le t^m$;\item if $t\le u$, then $u\le t^m$.\end{enumerate}
In particular, $t^m$ is the unique maximal element in $t^{\uparrow}$.
\begin{proof}
If $t$ is idempotent, (1) is trivial and (2) follows from $E$-unitarity.
Otherwise let $t=\langle\pi;t_1,\ldots,t_n\rangle$ and let $k$ be the greatest $i$ such that $t_k\notin E$ or $\pi k\ne k$.
Then $t^m=\langle\pi;(t_1)^m,\ldots,(t_k)^m\rangle$.
By coinductive hypothesis $t_i\le(t_i)^m$ for $1\le i\le k$, and $t_i\in E$ for $k<i\le n$, so that by Lemma 18 $t\le t^m$, showing (1).
Concerning (2), let $u=\langle\rho;u_1,\ldots,u_m\rangle$ such that $t\le u$.
Without loss of generality, we assume that $n=m$, so that $k\le m$, where $k$ is defined above.
Moreover, we have $\pi=\rho$, so that by Lemma 18 and by coinductive hypothesis, $u_i\le(t_i)^m$ for $1\le i\le k$.
As $t_i\le u_i$ and $t_i\in E$ for $k<i\le n$, we have $u_i\in E$ for $k<i\le m$ by $E$-unitarity.
By Lemma 18 we conclude that $u\le t^m$.
\end{proof}
We now prove that $\mathsf{PT}$ is $F$-inverse. The proof relies on the following general lemma that will be used multiple times.
\paragraph{Lemma 21.} Let $S$ be a $E$-unitary inverse semigroup. Assume that for every $t\in S$, the set $t^{\uparrow}$ has a unique maximal element $t^m$.
Then for all $t,u\in S$, $t\mathrel{\sigma_S}u$ iff $t^m=u^m$.
\begin{proof}
Assuming $t\mathrel{\sigma_S}u$, there exists $w$ such that $w\le t,u$. Hence $w^m=t^m=u^m$.
Conversely, let $w=t^m=u^m$. Then $u=we$ for some $e\in E(S)$ and $t=dw$ for some $d\in E(S)$.
Moreover, $t^*u=t^*we\le w^*we\in E(S)$ and $tu^*=dwu^*\le dww^*\in E(S)$, so that $t,u$ are compatible, i.e., $t\sim u$.
As $S$ is $E$-unitary, $\sim\,=\sigma_S$, so that $t\mathrel{\sigma_S}u$.
\end{proof}

\paragraph{Theorem 22.} The monoid $\mathsf{PT}$ is $F$-inverse.

\begin{proof}
By Lemma 21 and Proposition 20.
\end{proof}

